
\documentclass[11pt,a4paper]{article}

%========================
% PACKAGES
%========================
\usepackage[utf8]{inputenc}
\usepackage[T1]{fontenc}
\usepackage[french]{babel}
\usepackage[a4paper,margin=2cm]{geometry}
\usepackage{amsmath,amssymb}
\usepackage{lmodern}
\usepackage{fancyhdr}
\usepackage{lastpage}
\usepackage{enumitem}
\usepackage{array,tabularx}
\usepackage{tikz}
\usepackage{pgfplots}
\usepackage[most]{tcolorbox}
\usepackage{xcolor}
\usepackage{hyperref}

\pgfplotsset{compat=1.18}

%========================
% COULEURS
%========================
\definecolor{bleu}{RGB}{30,80,160}
\definecolor{vert}{RGB}{35,130,80}
\definecolor{rouge}{RGB}{180,45,45}
\definecolor{orange}{RGB}{220,130,30}
\definecolor{grisclair}{RGB}{245,245,245}

%========================
% MISE EN PAGE
%========================
\setlength{\headheight}{14pt}
\pagestyle{fancy}
\fancyhf{}
\fancyhead[L]{Première spécialité mathématiques}
\fancyhead[C]{TD -- Fonction exponentielle}
\fancyhead[R]{Page \thepage/\pageref{LastPage}}
\fancyfoot[C]{Pierre Garcia-Buscaylet -- TD complet}

\hypersetup{
    colorlinks=true,
    linkcolor=bleu,
    urlcolor=bleu
}

%========================
% COMMANDES
%========================
\newcommand{\R}{\mathbb{R}}
\newcommand{\e}{\mathrm{e}}
\newcommand{\etoiles}[1]{\textcolor{orange}{#1}}
\newcommand{\competences}[1]{\smallskip\noindent\textbf{Compétences travaillées :} #1\par\smallskip}

%========================
% BOÎTES
%========================
\tcbset{
  arc=2mm,
  boxrule=0.7pt,
  left=2mm,
  right=2mm,
  top=1mm,
  bottom=1mm
}

\newtcolorbox{rappelbox}{
  title=\textbf{Rappel essentiel},
  colback=blue!4,
  colframe=bleu
}

\newtcolorbox{methodebox}{
  title=\textbf{Méthode},
  colback=orange!8,
  colframe=orange
}

\newtcolorbox{attentionbox}{
  title=\textbf{Attention},
  colback=red!5,
  colframe=rouge
}

\newtcolorbox{exercicebox}[1]{
  title=\textbf{#1},
  colback=white,
  colframe=bleu
}

\newtcolorbox{correctionbox}[1]{
  title=\textbf{#1},
  colback=green!4,
  colframe=vert
}

%========================
% COMPTEURS
%========================
\newcounter{exercice}
\newenvironment{exercice}[1]{
\refstepcounter{exercice}
\begin{exercicebox}{Exercice \theexercice{} -- #1}
}{
\end{exercicebox}
}

\newcounter{correction}
\newenvironment{correction}[1]{
\refstepcounter{correction}
\begin{correctionbox}{Correction de l'exercice \thecorrection{} -- #1}
}{
\end{correctionbox}
}

%========================
% DOCUMENT
%========================
\begin{document}

%========================
% TITRE
%========================
\begin{titlepage}
\centering

\vspace*{2cm}

{\Huge \textbf{TD complet}}\\[0.3cm]
{\Huge \textbf{Fonction exponentielle}}\\[0.5cm]
{\Large Première générale -- spécialité mathématiques}\\[1cm]

\begin{tcolorbox}[colback=blue!5,colframe=bleu,width=0.85\textwidth]
\centering
\Large
Calculs, équations, inéquations, dérivation, variations, suites, modélisation et problèmes longs.
\end{tcolorbox}

\vfill

\begin{tabular}{|>{\bfseries}l|l|}
\hline
Niveau & Première spécialité mathématiques \\
\hline
Chapitre & Fonction exponentielle \\
\hline
Objectif & Maîtriser les méthodes fondamentales du chapitre \\
\hline
Format & Énoncés puis corrigé complet \\
\hline
\end{tabular}

\vfill

{\large Document prêt à compiler}

\end{titlepage}

\tableofcontents
\newpage

%================================================
\section{Rappels essentiels de cours}

\begin{rappelbox}
La fonction exponentielle est l'unique fonction dérivable sur $\R$ telle que :
\[
\exp'(x)=\exp(x)
\quad \text{et} \quad
\exp(0)=1.
\]
On note :
\[
\exp(x)=\e^x.
\]
Le nombre $\e$ est défini par :
\[
\e=\exp(1)\approx 2{,}718.
\]
\end{rappelbox}

\begin{rappelbox}
Pour tous réels $a$ et $b$ :
\[
\e^{a+b}=\e^a\e^b,
\qquad
\e^{-a}=\frac{1}{\e^a},
\qquad
\e^{a-b}=\frac{\e^a}{\e^b}.
\]
Pour tout entier relatif $n$ :
\[
\e^{na}=(\e^a)^n.
\]
\end{rappelbox}

\begin{rappelbox}
Pour tout réel $x$ :
\[
\e^x>0.
\]
De plus :
\[
(\e^x)'=\e^x.
\]
La fonction exponentielle est donc strictement croissante sur $\R$.
\end{rappelbox}

\begin{rappelbox}
Comme la fonction exponentielle est strictement croissante, pour tous réels $A$ et $B$ :
\[
\e^A=\e^B \Longleftrightarrow A=B,
\]
et :
\[
\e^A<\e^B \Longleftrightarrow A<B.
\]
\end{rappelbox}

\begin{attentionbox}
En Première, on évite d'utiliser le logarithme comme outil principal de résolution. On résout surtout les équations de la forme :
\[
\e^A=\e^B.
\]
\end{attentionbox}

%================================================
\section{Méthodes types}

\begin{methodebox}
\textbf{Méthode 1 -- Simplifier une expression exponentielle}

Pour simplifier une expression :
\[
\e^a\e^b,
\]
on additionne les exposants :
\[
\e^a\e^b=\e^{a+b}.
\]

Pour simplifier :
\[
\frac{\e^a}{\e^b},
\]
on soustrait les exposants :
\[
\frac{\e^a}{\e^b}=\e^{a-b}.
\]
\end{methodebox}

\begin{methodebox}
\textbf{Méthode 2 -- Résoudre une équation du type $\e^A=\e^B$}

On utilise l'injectivité de l'exponentielle :
\[
\e^A=\e^B \Longleftrightarrow A=B.
\]
On résout ensuite l'équation obtenue.
\end{methodebox}

\begin{methodebox}
\textbf{Méthode 3 -- Résoudre une inéquation du type $\e^A<\e^B$}

Comme l'exponentielle est strictement croissante :
\[
\e^A<\e^B \Longleftrightarrow A<B.
\]
Le sens de l'inégalité ne change pas.
\end{methodebox}

\begin{methodebox}
\textbf{Méthode 4 -- Dériver une fonction contenant $\e^x$}

On utilise :
\[
(\e^x)'=\e^x.
\]
Si la fonction est un produit, on utilise :
\[
(uv)'=u'v+uv'.
\]
\end{methodebox}

\begin{methodebox}
\textbf{Méthode 5 -- Étudier les variations}

Pour étudier les variations d'une fonction contenant $\e^x$ :
\begin{enumerate}
    \item calculer la dérivée ;
    \item factoriser si possible par $\e^x$ ;
    \item utiliser $\e^x>0$ ;
    \item étudier le signe du facteur restant ;
    \item construire le tableau de variations.
\end{enumerate}
\end{methodebox}

\begin{methodebox}
\textbf{Méthode 6 -- Modéliser par une exponentielle}

Une fonction du type :
\[
f(t)=C\e^{kt}
\]
modélise souvent une évolution exponentielle.
\begin{itemize}
    \item $C=f(0)$ est la valeur initiale.
    \item Si $k>0$, le phénomène est croissant.
    \item Si $k<0$, le phénomène est décroissant.
\end{itemize}
\end{methodebox}

\newpage

%================================================
\section{Exercices d'application directe}

\begin{exercice}{APP -- \etoiles{$\star$} -- Valeurs particulières}
\competences{Connaître les valeurs fondamentales de l'exponentielle.}
Calculer :
\begin{enumerate}[label=\alph*)]
    \item $\e^0$
    \item $\e^1$
    \item $\e^{-1}$
    \item $\e^2\e^{-2}$
    \item $\dfrac{\e^5}{\e^5}$
\end{enumerate}
\end{exercice}

\begin{exercice}{APP -- \etoiles{$\star$} -- Produits d'exponentielles}
\competences{Utiliser $\e^a\e^b=\e^{a+b}$.}
Simplifier :
\begin{enumerate}[label=\alph*)]
    \item $\e^2\e^5$
    \item $\e^3\e^{-1}$
    \item $\e^x\e^4$
    \item $\e^{x+1}\e^{2x-3}$
    \item $\e^{-x}\e^{3x+1}$
\end{enumerate}
\end{exercice}

\begin{exercice}{APP -- \etoiles{$\star$} -- Quotients d'exponentielles}
\competences{Utiliser $\dfrac{\e^a}{\e^b}=\e^{a-b}$.}
Simplifier :
\begin{enumerate}[label=\alph*)]
    \item $\dfrac{\e^7}{\e^2}$
    \item $\dfrac{\e^x}{\e^3}$
    \item $\dfrac{\e^{4x}}{\e^{x-2}}$
    \item $\dfrac{\e^{2x+1}}{\e^{x+5}}$
    \item $\dfrac{\e^{-x+3}}{\e^{2x-1}}$
\end{enumerate}
\end{exercice}

\begin{exercice}{APP -- \etoiles{$\star\star$} -- Écriture sous forme d'une seule exponentielle}
\competences{Regrouper produits et quotients d'exponentielles.}
Écrire sous la forme $\e^{A}$ :
\begin{enumerate}[label=\alph*)]
    \item $\e^{x+2}\e^{3x-1}$
    \item $\dfrac{\e^{5x+1}}{\e^{2x-4}}$
    \item $\dfrac{\e^{x+1}\e^{2x-3}}{\e^{4}}$
    \item $\dfrac{\e^{3x}\e^{-x+2}}{\e^{x-1}}$
\end{enumerate}
\end{exercice}

\begin{exercice}{APP -- \etoiles{$\star$} -- Signe}
\competences{Utiliser la positivité de l'exponentielle.}
Dire si les affirmations suivantes sont vraies ou fausses. Justifier.
\begin{enumerate}[label=\alph*)]
    \item Pour tout réel $x$, $\e^x>0$.
    \item Il existe un réel $x$ tel que $\e^x=0$.
    \item Pour tout réel $x$, $-3\e^x<0$.
    \item Pour tout réel $x$, $(x-1)\e^x$ a le même signe que $x-1$.
\end{enumerate}
\end{exercice}

\begin{exercice}{APP -- \etoiles{$\star\star$} -- Équations simples}
\competences{Résoudre $\e^A=\e^B$.}
Résoudre dans $\R$ :
\begin{enumerate}[label=\alph*)]
    \item $\e^x=\e^3$
    \item $\e^{x+1}=\e^5$
    \item $\e^{2x-1}=\e^{x+4}$
    \item $\e^{-3x}=\e^6$
\end{enumerate}
\end{exercice}

\begin{exercice}{APP -- \etoiles{$\star\star$} -- Inéquations simples}
\competences{Résoudre $\e^A<\e^B$ ou $\e^A\geq \e^B$.}
Résoudre dans $\R$ :
\begin{enumerate}[label=\alph*)]
    \item $\e^x<\e^2$
    \item $\e^{x-1}\leq \e^4$
    \item $\e^{2x+1}\geq \e^5$
    \item $\e^{-x}<\e^3$
\end{enumerate}
\end{exercice}

\begin{exercice}{APP -- \etoiles{$\star\star$} -- Dérivées directes}
\competences{Dériver des fonctions simples contenant $\e^x$.}
Dériver :
\begin{enumerate}[label=\alph*)]
    \item $f(x)=3\e^x-5$
    \item $g(x)=\e^x+2x-1$
    \item $h(x)=-4\e^x+7x$
    \item $p(x)=5\e^x-x^2$
\end{enumerate}
\end{exercice}

%================================================
\section{Exercices de méthode}

\begin{exercice}{METH -- \etoiles{$\star\star$} -- Factorisation par $\e^x$}
\competences{Factoriser une expression contenant $\e^x$.}
Factoriser :
\begin{enumerate}[label=\alph*)]
    \item $3\e^x+x\e^x$
    \item $(2x-1)\e^x+5\e^x$
    \item $x^2\e^x-4x\e^x$
    \item $(x+1)\e^x-(2x-3)\e^x$
\end{enumerate}
\end{exercice}

\begin{exercice}{METH -- \etoiles{$\star\star$} -- Produit avec $\e^x$}
\competences{Utiliser la formule du produit.}
Dériver :
\begin{enumerate}[label=\alph*)]
    \item $f(x)=x\e^x$
    \item $g(x)=(x+2)\e^x$
    \item $h(x)=(2x-3)\e^x$
    \item $p(x)=(x^2+1)\e^x$
\end{enumerate}
\end{exercice}

\begin{exercice}{METH -- \etoiles{$\star\star$} -- Signe d'une dérivée}
\competences{Utiliser $\e^x>0$ pour étudier un signe.}
Étudier le signe des expressions suivantes sur $\R$ :
\begin{enumerate}[label=\alph*)]
    \item $(x-2)\e^x$
    \item $(3x+6)\e^x$
    \item $-(x+1)\e^x$
    \item $(x^2-4)\e^x$
\end{enumerate}
\end{exercice}

\begin{exercice}{METH -- \etoiles{$\star\star\star$} -- Variations de $x\e^x$}
\competences{Dériver, étudier le signe, dresser un tableau de variations.}
On considère :
\[
f(x)=x\e^x.
\]
\begin{enumerate}
    \item Justifier que $f$ est définie sur $\R$.
    \item Calculer $f'(x)$.
    \item Factoriser $f'(x)$.
    \item Étudier le signe de $f'(x)$.
    \item Dresser le tableau de variations de $f$.
\end{enumerate}
\end{exercice}

\begin{exercice}{METH -- \etoiles{$\star\star\star$} -- Étude de $(x-1)\e^x$}
\competences{Faire une étude de variations complète.}
On considère :
\[
f(x)=(x-1)\e^x.
\]
\begin{enumerate}
    \item Calculer $f'(x)$.
    \item Montrer que $f'(x)=x\e^x$.
    \item Étudier le signe de $f'(x)$.
    \item En déduire les variations de $f$.
    \item Déterminer l'extremum éventuel de $f$.
\end{enumerate}
\end{exercice}

\begin{exercice}{METH -- \etoiles{$\star\star\star$} -- Fonction $\e^x-x$}
\competences{Étudier une fonction avec exponentielle et affine.}
On considère :
\[
f(x)=\e^x-x.
\]
\begin{enumerate}
    \item Calculer $f'(x)$.
    \item Résoudre $\e^x-1=0$.
    \item Étudier le signe de $f'(x)$.
    \item Dresser le tableau de variations de $f$.
    \item Déterminer le minimum de $f$.
\end{enumerate}
\end{exercice}

\begin{exercice}{METH -- \etoiles{$\star\star\star$} -- Tangente à $\e^x$}
\competences{Déterminer une équation de tangente.}
On note $C$ la courbe de la fonction $f(x)=\e^x$.
\begin{enumerate}
    \item Calculer $f(0)$ et $f'(0)$.
    \item En déduire l'équation de la tangente à $C$ au point d'abscisse $0$.
    \item Calculer $f(1)$ et $f'(1)$.
    \item Donner l'équation de la tangente à $C$ au point d'abscisse $1$.
\end{enumerate}
\end{exercice}

\begin{exercice}{METH -- \etoiles{$\star\star\star$} -- Suite exponentielle}
\competences{Relier exponentielle et suite géométrique.}
On définit :
\[
u_n=\e^{2n+1}.
\]
\begin{enumerate}
    \item Calculer $u_0$ et $u_1$.
    \item Montrer que $(u_n)$ est une suite géométrique.
    \item Donner sa raison.
    \item Déterminer son sens de variation.
\end{enumerate}
\end{exercice}

%================================================
\section{Exercices de démonstration et raisonnement}

\begin{exercice}{DEM -- \etoiles{$\star\star$} -- Propriété de l'inverse}
\competences{Démontrer $\e^{-a}=1/\e^a$.}
Soit $a$ un réel.
\begin{enumerate}
    \item Écrire $\e^{a+(-a)}$ de deux façons.
    \item En déduire que $\e^a\e^{-a}=1$.
    \item Conclure que :
    \[
    \e^{-a}=\frac{1}{\e^a}.
    \]
\end{enumerate}
\end{exercice}

\begin{exercice}{DEM -- \etoiles{$\star\star$} -- Propriété du quotient}
\competences{Démontrer $\e^{a-b}=\e^a/\e^b$.}
Soient $a$ et $b$ deux réels.
\begin{enumerate}
    \item Écrire $a-b$ sous la forme $a+(-b)$.
    \item Utiliser la propriété $\e^{u+v}=\e^u\e^v$.
    \item Utiliser la propriété de l'inverse.
    \item Conclure.
\end{enumerate}
\end{exercice}

\begin{exercice}{DEM -- \etoiles{$\star\star\star$} -- Croissance et équations}
\competences{Justifier une résolution par injectivité.}
On suppose que $A$ et $B$ sont deux expressions réelles.
\begin{enumerate}
    \item Rappeler pourquoi la fonction exponentielle est strictement croissante.
    \item Expliquer pourquoi une fonction strictement croissante est injective.
    \item Justifier alors :
    \[
    \e^A=\e^B \Longleftrightarrow A=B.
    \]
\end{enumerate}
\end{exercice}

\begin{exercice}{DEM -- \etoiles{$\star\star\star$} -- Signe d'une dérivée factorisée}
\competences{Rédiger un raisonnement de signe.}
Soit :
\[
f'(x)=(2x-6)\e^x.
\]
\begin{enumerate}
    \item Justifier que $\e^x>0$ pour tout réel $x$.
    \item Expliquer pourquoi le signe de $f'(x)$ est celui de $2x-6$.
    \item En déduire le signe de $f'(x)$ selon les valeurs de $x$.
\end{enumerate}
\end{exercice}

\begin{exercice}{DEM -- \etoiles{$\star\star\star$} -- Suite géométrique générale}
\competences{Démontrer qu'une suite exponentielle est géométrique.}
Soient $a$ et $b$ deux réels. On définit :
\[
u_n=\e^{an+b}.
\]
\begin{enumerate}
    \item Calculer $u_{n+1}$.
    \item Calculer $\dfrac{u_{n+1}}{u_n}$.
    \item En déduire que $(u_n)$ est géométrique.
    \item Donner sa raison.
\end{enumerate}
\end{exercice}

\begin{exercice}{DEM -- \etoiles{$\star\star\star$} -- Minimum de $\e^x-x$}
\competences{Démontrer une inégalité à l'aide d'une étude de fonction.}
On considère :
\[
f(x)=\e^x-x.
\]
\begin{enumerate}
    \item Étudier les variations de $f$.
    \item Montrer que $f$ admet un minimum.
    \item En déduire que, pour tout réel $x$ :
    \[
    \e^x\geq x.
    \]
\end{enumerate}
\end{exercice}

%================================================
\section{Exercices de réflexion}

\begin{exercice}{REF -- \etoiles{$\star\star\star$} -- Équation du second degré}
\competences{Résoudre une équation exponentielle se ramenant au second degré.}
Résoudre dans $\R$ :
\[
\e^{2x}-5\e^x+6=0.
\]
On pourra poser $X=\e^x$.
\end{exercice}

\begin{exercice}{REF -- \etoiles{$\star\star\star$} -- Inéquation avec trinôme}
\competences{Transformer une inéquation exponentielle.}
Résoudre dans $\R$ :
\[
\e^{x^2-3x}\leq 1.
\]
\end{exercice}

\begin{exercice}{REF -- \etoiles{$\star\star\star$} -- Position relative}
\competences{Comparer deux fonctions à l'aide d'une différence.}
On considère :
\[
f(x)=\e^x
\quad \text{et} \quad
g(x)=x+1.
\]
\begin{enumerate}
    \item Étudier la fonction $h(x)=\e^x-x-1$ dans un cadre accessible.
    \item Montrer que $h(0)=0$.
    \item À l'aide de l'étude de $h$, conjecturer puis justifier que la courbe de $\e^x$ est au-dessus de sa tangente en $0$.
\end{enumerate}
\end{exercice}

\begin{exercice}{REF -- \etoiles{$\star\star\star\star$} -- Étude d'une fonction}
\competences{Étude complète avec dérivation et variations.}
On considère :
\[
f(x)=(x+2)\e^x-3.
\]
\begin{enumerate}
    \item Calculer $f'(x)$.
    \item Factoriser $f'(x)$.
    \item Étudier le signe de $f'(x)$.
    \item Dresser le tableau de variations de $f$.
    \item Déterminer l'éventuel extremum de $f$.
    \item Résoudre $f(x)=0$ lorsque cela est possible exactement ou expliquer pourquoi une valeur approchée peut être nécessaire.
\end{enumerate}
\end{exercice}

\begin{exercice}{REF -- \etoiles{$\star\star\star\star$} -- Paramètre}
\competences{Raisonner avec un paramètre simple.}
Soit $a$ un réel. On considère :
\[
f_a(x)=\e^x-ax.
\]
\begin{enumerate}
    \item Calculer $f_a'(x)$.
    \item Dans le cas $a=1$, retrouver les variations de $\e^x-x$.
    \item Dans le cas $a=-2$, étudier le signe de $f_a'(x)$.
    \item Discuter qualitativement l'influence du signe de $a$.
\end{enumerate}
\end{exercice}

\begin{exercice}{REF -- \etoiles{$\star\star\star\star$} -- Tangente à une fonction produit}
\competences{Calculer une tangente à une fonction contenant $\e^x$.}
On considère :
\[
f(x)=(x-1)\e^x.
\]
\begin{enumerate}
    \item Calculer $f(0)$.
    \item Calculer $f'(x)$ puis $f'(0)$.
    \item Donner l'équation de la tangente à la courbe de $f$ au point d'abscisse $0$.
    \item Même question au point d'abscisse $1$.
\end{enumerate}
\end{exercice}

%================================================
\section{Exercices de modélisation}

\begin{exercice}{MOD -- \etoiles{$\star\star$} -- Population}
\competences{Interpréter un modèle exponentiel.}
Une population est modélisée par :
\[
P(t)=1200\e^{0,04t},
\]
où $t$ est exprimé en années.
\begin{enumerate}
    \item Calculer $P(0)$.
    \item Le phénomène est-il croissant ou décroissant ?
    \item Calculer une valeur approchée de $P(5)$.
    \item Interpréter le coefficient $0,04$.
\end{enumerate}
\end{exercice}

\begin{exercice}{MOD -- \etoiles{$\star\star$} -- Médicament}
\competences{Modéliser une décroissance exponentielle.}
La quantité d'un médicament dans le sang est donnée par :
\[
M(t)=80\e^{-0,3t},
\]
où $t$ est exprimé en heures.
\begin{enumerate}
    \item Quelle est la quantité initiale ?
    \item Le médicament augmente-t-il ou diminue-t-il dans le sang ?
    \item Calculer $M(4)$ à $0,1$ près.
    \item Interpréter le résultat.
\end{enumerate}
\end{exercice}

\begin{exercice}{MOD -- \etoiles{$\star\star$} -- Radioactivité}
\competences{Lire un modèle exponentiel décroissant.}
Une substance radioactive est modélisée par :
\[
Q(t)=500\e^{-0,12t}.
\]
\begin{enumerate}
    \item Calculer $Q(0)$.
    \item Justifier que la quantité diminue.
    \item Calculer $Q(10)$ à l'unité près.
    \item Que représente le nombre $500$ ?
\end{enumerate}
\end{exercice}

\begin{exercice}{MOD -- \etoiles{$\star\star\star$} -- Capital}
\competences{Calculer et interpréter une croissance exponentielle.}
Un capital est modélisé par :
\[
C(t)=2000\e^{0,025t}.
\]
\begin{enumerate}
    \item Calculer le capital initial.
    \item Calculer le capital après $8$ ans.
    \item Le modèle traduit-il une croissance ou une décroissance ?
    \item Comparer avec une croissance linéaire de $50$ euros par an.
\end{enumerate}
\end{exercice}

\begin{exercice}{MOD -- \etoiles{$\star\star\star$} -- Dilution}
\competences{Interpréter une décroissance dans un contexte chimique.}
La concentration d'une solution est modélisée par :
\[
c(t)=15\e^{-0,2t}.
\]
\begin{enumerate}
    \item Donner la concentration initiale.
    \item Calculer $c(3)$.
    \item Justifier que la concentration diminue.
    \item Interpréter le coefficient $-0,2$.
\end{enumerate}
\end{exercice}

\begin{exercice}{MOD -- \etoiles{$\star\star\star$} -- Refroidissement simplifié}
\competences{Exploiter un modèle de décroissance.}
La différence de température entre un objet et l'air ambiant est modélisée par :
\[
D(t)=70\e^{-0,15t}.
\]
\begin{enumerate}
    \item Calculer $D(0)$.
    \item Calculer $D(10)$.
    \item La différence de température augmente-t-elle ou diminue-t-elle ?
    \item Donner une interprétation concrète du résultat.
\end{enumerate}
\end{exercice}

%================================================
\section{Problèmes longs type contrôle}

\begin{exercice}{PROB -- \etoiles{$\star\star\star\star$} -- Étude complète}
\competences{Dériver, factoriser, étudier des variations et interpréter.}
On considère la fonction :
\[
f(x)=(x+1)\e^x-2.
\]
\begin{enumerate}
    \item Montrer que $f$ est définie sur $\R$.
    \item Calculer $f'(x)$.
    \item Montrer que :
    \[
    f'(x)=(x+2)\e^x.
    \]
    \item Étudier le signe de $f'(x)$.
    \item Dresser le tableau de variations de $f$.
    \item Calculer $f(-2)$.
    \item En déduire l'existence d'un extremum et préciser sa nature.
\end{enumerate}
\end{exercice}

\begin{exercice}{PROB -- \etoiles{$\star\star\star\star$} -- Suites et exponentielle}
\competences{Relier suite géométrique et exponentielle.}
On définit, pour tout entier naturel $n$ :
\[
u_n=4\e^{-0,5n+2}.
\]
\begin{enumerate}
    \item Calculer $u_0$ et $u_1$.
    \item Montrer que $(u_n)$ est une suite géométrique.
    \item Donner sa raison.
    \item La suite est-elle croissante ou décroissante ?
    \item Interpréter cette suite comme un modèle de décroissance.
\end{enumerate}
\end{exercice}

\begin{exercice}{PROB -- \etoiles{$\star\star\star\star\star$} -- Modélisation et dérivation}
\competences{Étudier un modèle, calculer une dérivée et interpréter.}
Une population de bactéries est modélisée par :
\[
N(t)=1000\e^{0,6t},
\]
où $t$ est exprimé en heures.
\begin{enumerate}
    \item Calculer $N(0)$.
    \item Justifier que la population augmente.
    \item Calculer $N(3)$ à l'unité près.
    \item Calculer $N'(t)$.
    \item Interpréter $N'(t)$ dans le contexte.
    \item Calculer $N'(0)$ et interpréter.
\end{enumerate}
\end{exercice}

\begin{exercice}{PROB -- \etoiles{$\star\star\star\star\star$} -- Fonction, tangente et équation}
\competences{Mélanger dérivation, tangente, équation et signe.}
On considère :
\[
f(x)=\e^x-2x.
\]
\begin{enumerate}
    \item Calculer $f'(x)$.
    \item Résoudre $f'(x)=0$ lorsque cela est possible sans logarithme.
    \item Étudier le signe de $f'(x)$ en utilisant la croissance de $\e^x$.
    \item Dresser le tableau de variations de $f$.
    \item Donner l'équation de la tangente à la courbe de $f$ au point d'abscisse $0$.
    \item Résoudre l'équation $f(x)=1$ si possible exactement.
\end{enumerate}
\end{exercice}

\newpage

%================================================
\section{Corrigé complet}

\begin{correction}{APP -- Valeurs particulières}
\begin{enumerate}[label=\alph*)]
    \item Par définition, $\e^0=1$.
    \item $\e^1=\e$.
    \item $\e^{-1}=\dfrac{1}{\e}$.
    \item $\e^2\e^{-2}=\e^{2-2}=\e^0=1$.
    \item $\dfrac{\e^5}{\e^5}=\e^{5-5}=\e^0=1$.
\end{enumerate}
\end{correction}

\begin{correction}{APP -- Produits d'exponentielles}
On utilise systématiquement :
\[
\e^a\e^b=\e^{a+b}.
\]
\begin{enumerate}[label=\alph*)]
    \item $\e^2\e^5=\e^{7}$.
    \item $\e^3\e^{-1}=\e^{2}$.
    \item $\e^x\e^4=\e^{x+4}$.
    \item $\e^{x+1}\e^{2x-3}=\e^{3x-2}$.
    \item $\e^{-x}\e^{3x+1}=\e^{2x+1}$.
\end{enumerate}
\end{correction}

\begin{correction}{APP -- Quotients d'exponentielles}
On utilise :
\[
\frac{\e^a}{\e^b}=\e^{a-b}.
\]
\begin{enumerate}[label=\alph*)]
    \item $\dfrac{\e^7}{\e^2}=\e^5$.
    \item $\dfrac{\e^x}{\e^3}=\e^{x-3}$.
    \item $\dfrac{\e^{4x}}{\e^{x-2}}=\e^{4x-(x-2)}=\e^{3x+2}$.
    \item $\dfrac{\e^{2x+1}}{\e^{x+5}}=\e^{x-4}$.
    \item $\dfrac{\e^{-x+3}}{\e^{2x-1}}=\e^{-x+3-2x+1}=\e^{-3x+4}$.
\end{enumerate}
\end{correction}

\begin{correction}{APP -- Écriture sous forme d'une seule exponentielle}
\begin{enumerate}[label=\alph*)]
    \item $\e^{x+2}\e^{3x-1}=\e^{4x+1}$.
    \item $\dfrac{\e^{5x+1}}{\e^{2x-4}}=\e^{3x+5}$.
    \item $\dfrac{\e^{x+1}\e^{2x-3}}{\e^4}=\dfrac{\e^{3x-2}}{\e^4}=\e^{3x-6}$.
    \item $\dfrac{\e^{3x}\e^{-x+2}}{\e^{x-1}}=\dfrac{\e^{2x+2}}{\e^{x-1}}=\e^{x+3}$.
\end{enumerate}
\end{correction}

\begin{correction}{APP -- Signe}
\begin{enumerate}[label=\alph*)]
    \item Vrai : pour tout réel $x$, $\e^x>0$.
    \item Faux : l'exponentielle ne s'annule jamais.
    \item Vrai : $\e^x>0$, donc $-3\e^x<0$.
    \item Vrai : comme $\e^x>0$, multiplier par $\e^x$ ne change pas le signe.
\end{enumerate}
\end{correction}

\begin{correction}{APP -- Équations simples}
\begin{enumerate}[label=\alph*)]
    \item $\e^x=\e^3 \Longleftrightarrow x=3$.
    \item $\e^{x+1}=\e^5 \Longleftrightarrow x+1=5$, donc $x=4$.
    \item $\e^{2x-1}=\e^{x+4}\Longleftrightarrow 2x-1=x+4$, donc $x=5$.
    \item $\e^{-3x}=\e^6\Longleftrightarrow -3x=6$, donc $x=-2$.
\end{enumerate}
\end{correction}

\begin{correction}{APP -- Inéquations simples}
\begin{enumerate}[label=\alph*)]
    \item $\e^x<\e^2\Longleftrightarrow x<2$.
    \item $\e^{x-1}\leq \e^4\Longleftrightarrow x-1\leq4$, donc $x\leq5$.
    \item $\e^{2x+1}\geq\e^5\Longleftrightarrow 2x+1\geq5$, donc $x\geq2$.
    \item $\e^{-x}<\e^3\Longleftrightarrow -x<3$, donc $x>-3$.
\end{enumerate}
\end{correction}

\begin{correction}{APP -- Dérivées directes}
\begin{enumerate}[label=\alph*)]
    \item $f'(x)=3\e^x$.
    \item $g'(x)=\e^x+2$.
    \item $h'(x)=-4\e^x+7$.
    \item $p'(x)=5\e^x-2x$.
\end{enumerate}
\end{correction}

\begin{correction}{METH -- Factorisation par $\e^x$}
\begin{enumerate}[label=\alph*)]
    \item $3\e^x+x\e^x=(x+3)\e^x$.
    \item $(2x-1)\e^x+5\e^x=(2x+4)\e^x$.
    \item $x^2\e^x-4x\e^x=x(x-4)\e^x$.
    \item $(x+1)\e^x-(2x-3)\e^x=(-x+4)\e^x$.
\end{enumerate}
\end{correction}

\begin{correction}{METH -- Produit avec $\e^x$}
On utilise :
\[
(uv)'=u'v+uv'.
\]
\begin{enumerate}[label=\alph*)]
    \item $f'(x)=\e^x+x\e^x=(x+1)\e^x$.
    \item $g'(x)=\e^x+(x+2)\e^x=(x+3)\e^x$.
    \item $h'(x)=2\e^x+(2x-3)\e^x=(2x-1)\e^x$.
    \item $p'(x)=2x\e^x+(x^2+1)\e^x=(x^2+2x+1)\e^x=(x+1)^2\e^x$.
\end{enumerate}
\end{correction}

\begin{correction}{METH -- Signe d'une dérivée}
Comme $\e^x>0$, le signe est celui du facteur polynomial.
\begin{enumerate}[label=\alph*)]
    \item $(x-2)\e^x$ est négatif si $x<2$, nul si $x=2$, positif si $x>2$.
    \item $(3x+6)\e^x$ est du signe de $3(x+2)$ : négatif si $x<-2$, nul si $x=-2$, positif si $x>-2$.
    \item $-(x+1)\e^x$ est positif si $x<-1$, nul si $x=-1$, négatif si $x>-1$.
    \item $(x^2-4)\e^x$ est du signe de $(x-2)(x+2)$ : positif sur $]-\infty;-2]\cup[2;+\infty[$ et négatif sur $]-2;2[$.
\end{enumerate}
\end{correction}

\begin{correction}{METH -- Variations de $x\e^x$}
La fonction $f(x)=x\e^x$ est définie sur $\R$.

Par produit :
\[
f'(x)=1\cdot \e^x+x\e^x=(x+1)\e^x.
\]
Comme $\e^x>0$, le signe de $f'(x)$ est celui de $x+1$.

Donc $f'(x)<0$ si $x<-1$, $f'(-1)=0$, et $f'(x)>0$ si $x>-1$.

Ainsi, $f$ est décroissante sur $]-\infty;-1]$ puis croissante sur $[-1;+\infty[$.

Elle admet un minimum en $x=-1$ :
\[
f(-1)=-\e^{-1}=-\frac{1}{\e}.
\]
\end{correction}

\begin{correction}{METH -- Étude de $(x-1)\e^x$}
\[
f(x)=(x-1)\e^x.
\]
Par produit :
\[
f'(x)=1\cdot \e^x+(x-1)\e^x=x\e^x.
\]
Comme $\e^x>0$, le signe de $f'(x)$ est celui de $x$.

Donc $f$ est décroissante sur $]-\infty;0]$, puis croissante sur $[0;+\infty[$.

L'extremum est un minimum atteint en $0$ :
\[
f(0)=(0-1)\e^0=-1.
\]
\end{correction}

\begin{correction}{METH -- Fonction $\e^x-x$}
\[
f'(x)=\e^x-1.
\]
On résout :
\[
\e^x-1=0 \Longleftrightarrow \e^x=1=\e^0 \Longleftrightarrow x=0.
\]
Comme l'exponentielle est croissante :
\[
\e^x<1 \Longleftrightarrow x<0,
\]
et :
\[
\e^x>1 \Longleftrightarrow x>0.
\]
Donc $f$ est décroissante sur $]-\infty;0]$, puis croissante sur $[0;+\infty[$.

Le minimum est :
\[
f(0)=\e^0-0=1.
\]
\end{correction}

\begin{correction}{METH -- Tangente à $\e^x$}
Pour $f(x)=\e^x$, on a :
\[
f'(x)=\e^x.
\]
En $0$ :
\[
f(0)=1,\qquad f'(0)=1.
\]
La tangente en $0$ a pour équation :
\[
y=f'(0)(x-0)+f(0)=x+1.
\]
En $1$ :
\[
f(1)=\e,\qquad f'(1)=\e.
\]
La tangente en $1$ est :
\[
y=\e(x-1)+\e=\e x.
\]
\end{correction}

\begin{correction}{METH -- Suite exponentielle}
\[
u_n=\e^{2n+1}.
\]
On a :
\[
u_0=\e^1=\e,
\qquad
u_1=\e^3.
\]
Ensuite :
\[
\frac{u_{n+1}}{u_n}
=
\frac{\e^{2(n+1)+1}}{\e^{2n+1}}
=
\frac{\e^{2n+3}}{\e^{2n+1}}
=
\e^2.
\]
Donc $(u_n)$ est géométrique de raison $\e^2$.

Comme $\e^2>1$, la suite est croissante.
\end{correction}

\begin{correction}{DEM -- Propriété de l'inverse}
On sait que :
\[
a+(-a)=0.
\]
Donc :
\[
\e^{a+(-a)}=\e^0=1.
\]
Mais, par la propriété fondamentale :
\[
\e^{a+(-a)}=\e^a\e^{-a}.
\]
Ainsi :
\[
\e^a\e^{-a}=1.
\]
Comme $\e^a>0$, on peut diviser par $\e^a$ :
\[
\e^{-a}=\frac{1}{\e^a}.
\]
\end{correction}

\begin{correction}{DEM -- Propriété du quotient}
On écrit :
\[
a-b=a+(-b).
\]
Donc :
\[
\e^{a-b}=\e^{a+(-b)}=\e^a\e^{-b}.
\]
Or :
\[
\e^{-b}=\frac{1}{\e^b}.
\]
Donc :
\[
\e^{a-b}=\e^a\frac{1}{\e^b}=\frac{\e^a}{\e^b}.
\]
\end{correction}

\begin{correction}{DEM -- Croissance et équations}
La fonction exponentielle est dérivable et :
\[
(\e^x)'=\e^x.
\]
Or :
\[
\e^x>0
\]
pour tout réel $x$. Sa dérivée est donc strictement positive, donc la fonction exponentielle est strictement croissante.

Une fonction strictement croissante ne peut pas prendre deux fois la même valeur : elle est injective. Ainsi :
\[
\e^A=\e^B \Longleftrightarrow A=B.
\]
\end{correction}

\begin{correction}{DEM -- Signe d'une dérivée factorisée}
Pour tout réel $x$, $\e^x>0$.

Donc le signe de :
\[
f'(x)=(2x-6)\e^x
\]
est le signe de $2x-6$.

On résout :
\[
2x-6=0 \Longleftrightarrow x=3.
\]
Donc :
\[
f'(x)<0 \text{ si } x<3,
\]
\[
f'(3)=0,
\]
\[
f'(x)>0 \text{ si } x>3.
\]
\end{correction}

\begin{correction}{DEM -- Suite géométrique générale}
On a :
\[
u_n=\e^{an+b}.
\]
Alors :
\[
u_{n+1}=\e^{a(n+1)+b}=\e^{an+a+b}.
\]
Donc :
\[
\frac{u_{n+1}}{u_n}
=
\frac{\e^{an+a+b}}{\e^{an+b}}
=
\e^a.
\]
Le quotient $\dfrac{u_{n+1}}{u_n}$ est constant, donc $(u_n)$ est géométrique de raison :
\[
q=\e^a.
\]
\end{correction}

\begin{correction}{DEM -- Minimum de $\e^x-x$}
On considère :
\[
f(x)=\e^x-x.
\]
Alors :
\[
f'(x)=\e^x-1.
\]
On sait que :
\[
\e^x-1<0 \Longleftrightarrow x<0,
\]
\[
\e^x-1=0 \Longleftrightarrow x=0,
\]
\[
\e^x-1>0 \Longleftrightarrow x>0.
\]
Donc $f$ décroît sur $]-\infty;0]$ puis croît sur $[0;+\infty[$.

Elle admet donc un minimum en $0$ :
\[
f(0)=1.
\]
Donc pour tout réel $x$ :
\[
f(x)\geq 1.
\]
Ainsi :
\[
\e^x-x\geq 1,
\]
donc :
\[
\e^x\geq x+1.
\]
En particulier :
\[
\e^x\geq x.
\]
\end{correction}

\begin{correction}{REF -- Équation du second degré}
On pose :
\[
X=\e^x.
\]
Comme $\e^x>0$, on a $X>0$.

L'équation devient :
\[
X^2-5X+6=0.
\]
On factorise :
\[
X^2-5X+6=(X-2)(X-3).
\]
Donc :
\[
X=2 \quad \text{ou} \quad X=3.
\]
Ainsi :
\[
\e^x=2
\quad \text{ou} \quad
\e^x=3.
\]
Ces solutions existent et sont uniques car l'exponentielle est strictement croissante, mais elles ne s'expriment pas exactement sans logarithme. On peut les approximer à la calculatrice.
\end{correction}

\begin{correction}{REF -- Inéquation avec trinôme}
\[
\e^{x^2-3x}\leq 1.
\]
Comme $1=\e^0$, on obtient :
\[
\e^{x^2-3x}\leq \e^0.
\]
L'exponentielle étant croissante :
\[
x^2-3x\leq 0.
\]
On factorise :
\[
x(x-3)\leq 0.
\]
Le produit est négatif ou nul entre les racines :
\[
S=[0;3].
\]
\end{correction}

\begin{correction}{REF -- Position relative}
On pose :
\[
h(x)=\e^x-x-1.
\]
Alors :
\[
h'(x)=\e^x-1.
\]
Comme précédemment, $h'(x)<0$ si $x<0$, $h'(0)=0$, et $h'(x)>0$ si $x>0$.

Donc $h$ admet un minimum en $0$.

Or :
\[
h(0)=\e^0-0-1=0.
\]
Donc, pour tout réel $x$ :
\[
h(x)\geq 0.
\]
Ainsi :
\[
\e^x-x-1\geq 0,
\]
c'est-à-dire :
\[
\e^x\geq x+1.
\]
La courbe de $\e^x$ est donc au-dessus de la droite $y=x+1$, qui est sa tangente en $0$.
\end{correction}

\begin{correction}{REF -- Étude d'une fonction}
\[
f(x)=(x+2)\e^x-3.
\]
Par produit :
\[
f'(x)=\e^x+(x+2)\e^x=(x+3)\e^x.
\]
Comme $\e^x>0$, le signe de $f'(x)$ est celui de $x+3$.

Donc $f$ est décroissante sur $]-\infty;-3]$ puis croissante sur $[-3;+\infty[$.

L'extremum est un minimum atteint en $x=-3$ :
\[
f(-3)=(-1)\e^{-3}-3=-\e^{-3}-3.
\]
Pour résoudre $f(x)=0$, on doit résoudre :
\[
(x+2)\e^x=3.
\]
Cette équation ne se ramène pas simplement à $\e^A=\e^B$ ; une valeur approchée peut donc être nécessaire.
\end{correction}

\begin{correction}{REF -- Paramètre}
\[
f_a(x)=\e^x-ax.
\]
Alors :
\[
f_a'(x)=\e^x-a.
\]
Si $a=1$ :
\[
f_1'(x)=\e^x-1.
\]
On retrouve que $f_1$ décroît sur $]-\infty;0]$ puis croît sur $[0;+\infty[$.

Si $a=-2$ :
\[
f_{-2}'(x)=\e^x+2.
\]
Comme $\e^x>0$, on a :
\[
f_{-2}'(x)>0.
\]
Donc $f_{-2}$ est strictement croissante sur $\R$.

Qualitativement, si $a<0$, alors $\e^x-a=\e^x+|a|>0$, donc la fonction est croissante. Si $a>0$, le signe de $\e^x-a$ dépend de la comparaison entre $\e^x$ et $a$.
\end{correction}

\begin{correction}{REF -- Tangente à une fonction produit}
\[
f(x)=(x-1)\e^x.
\]
On a :
\[
f(0)=-1.
\]
La dérivée est :
\[
f'(x)=\e^x+(x-1)\e^x=x\e^x.
\]
Donc :
\[
f'(0)=0.
\]
La tangente en $0$ est :
\[
y=f'(0)(x-0)+f(0)=-1.
\]
En $1$ :
\[
f(1)=0,
\qquad
f'(1)=\e.
\]
La tangente en $1$ est :
\[
y=\e(x-1).
\]
\end{correction}

\begin{correction}{MOD -- Population}
\[
P(t)=1200\e^{0,04t}.
\]
\begin{enumerate}
    \item $P(0)=1200\e^0=1200$.
    \item Comme $0,04>0$, le phénomène est croissant.
    \item $P(5)=1200\e^{0,2}\approx 1200\times1,221=1465$.
    \item Le coefficient $0,04$ mesure le rythme de croissance exponentielle.
\end{enumerate}
\end{correction}

\begin{correction}{MOD -- Médicament}
\[
M(t)=80\e^{-0,3t}.
\]
\begin{enumerate}
    \item $M(0)=80$.
    \item Comme $-0,3<0$, la quantité diminue.
    \item $M(4)=80\e^{-1,2}\approx 80\times0,301=24,1$.
    \item Après $4$ heures, il reste environ $24,1$ unités de médicament dans le sang.
\end{enumerate}
\end{correction}

\begin{correction}{MOD -- Radioactivité}
\[
Q(t)=500\e^{-0,12t}.
\]
\begin{enumerate}
    \item $Q(0)=500$.
    \item Le coefficient $-0,12$ est négatif, donc la quantité diminue.
    \item $Q(10)=500\e^{-1,2}\approx 151$.
    \item $500$ représente la quantité initiale.
\end{enumerate}
\end{correction}

\begin{correction}{MOD -- Capital}
\[
C(t)=2000\e^{0,025t}.
\]
\begin{enumerate}
    \item $C(0)=2000$.
    \item $C(8)=2000\e^{0,2}\approx 2443$.
    \item Le phénomène est croissant car $0,025>0$.
    \item Une croissance linéaire de $50$ euros par an donnerait après $8$ ans :
    \[
    2000+8\times50=2400.
    \]
    Le modèle exponentiel donne ici environ $2443$, donc un peu plus.
\end{enumerate}
\end{correction}

\begin{correction}{MOD -- Dilution}
\[
c(t)=15\e^{-0,2t}.
\]
\begin{enumerate}
    \item $c(0)=15$.
    \item $c(3)=15\e^{-0,6}\approx 15\times0,549=8,24$.
    \item Comme $-0,2<0$, la concentration diminue.
    \item Le coefficient $-0,2$ traduit une décroissance exponentielle.
\end{enumerate}
\end{correction}

\begin{correction}{MOD -- Refroidissement simplifié}
\[
D(t)=70\e^{-0,15t}.
\]
\begin{enumerate}
    \item $D(0)=70$.
    \item $D(10)=70\e^{-1,5}\approx 70\times0,223=15,6$.
    \item La différence diminue car $-0,15<0$.
    \item Après $10$ unités de temps, l'écart avec l'air ambiant est d'environ $15,6$ degrés.
\end{enumerate}
\end{correction}

\begin{correction}{PROB -- Étude complète}
\[
f(x)=(x+1)\e^x-2.
\]
La fonction est définie sur $\R$.

Par produit :
\[
f'(x)=\e^x+(x+1)\e^x=(x+2)\e^x.
\]
Comme $\e^x>0$, le signe de $f'(x)$ est celui de $x+2$.

Donc :
\[
f'(x)<0 \text{ si } x<-2,
\]
\[
f'(-2)=0,
\]
\[
f'(x)>0 \text{ si } x>-2.
\]
La fonction est décroissante sur $]-\infty;-2]$ puis croissante sur $[-2;+\infty[$.

On calcule :
\[
f(-2)=(-1)\e^{-2}-2=-\e^{-2}-2.
\]
La fonction admet donc un minimum en $x=-2$, égal à :
\[
-\e^{-2}-2.
\]
\end{correction}

\begin{correction}{PROB -- Suites et exponentielle}
\[
u_n=4\e^{-0,5n+2}.
\]
On a :
\[
u_0=4\e^2,
\qquad
u_1=4\e^{1,5}.
\]
Puis :
\[
\frac{u_{n+1}}{u_n}
=
\frac{4\e^{-0,5(n+1)+2}}{4\e^{-0,5n+2}}
=
\e^{-0,5}.
\]
Donc $(u_n)$ est géométrique de raison :
\[
q=\e^{-0,5}.
\]
Comme :
\[
\e^{-0,5}<1,
\]
la suite est décroissante.

Elle peut modéliser une quantité qui diminue de manière multiplicative à chaque étape.
\end{correction}

\begin{correction}{PROB -- Modélisation et dérivation}
\[
N(t)=1000\e^{0,6t}.
\]
\begin{enumerate}
    \item $N(0)=1000$.
    \item Comme $0,6>0$, la population augmente.
    \item $N(3)=1000\e^{1,8}\approx 6049$.
    \item Dans ce modèle, on admet la dérivée :
    \[
    N'(t)=1000\times0,6\e^{0,6t}=600\e^{0,6t}.
    \]
    \item $N'(t)$ représente la vitesse instantanée de croissance de la population.
    \item $N'(0)=600$. Au départ, la population augmente à une vitesse d'environ $600$ bactéries par heure.
\end{enumerate}
\end{correction}

\begin{correction}{PROB -- Fonction, tangente et équation}
\[
f(x)=\e^x-2x.
\]
\[
f'(x)=\e^x-2.
\]
L'équation :
\[
f'(x)=0
\]
donne :
\[
\e^x=2.
\]
Cette équation admet une unique solution, car l'exponentielle est strictement croissante, mais elle ne s'exprime pas exactement sans logarithme.

On peut néanmoins raisonner ainsi :
\[
f'(x)<0 \Longleftrightarrow \e^x<2,
\]
\[
f'(x)>0 \Longleftrightarrow \e^x>2.
\]
Il existe donc un unique réel $\alpha$ tel que :
\[
\e^\alpha=2.
\]
Alors $f$ est décroissante sur $]-\infty;\alpha]$ puis croissante sur $[\alpha;+\infty[$.

Pour la tangente en $0$ :
\[
f(0)=1,
\qquad
f'(0)=1-2=-1.
\]
Donc la tangente est :
\[
y=-x+1.
\]

Enfin :
\[
f(x)=1
\Longleftrightarrow
\e^x-2x=1.
\]
On remarque que $x=0$ est solution :
\[
\e^0-0=1.
\]
Une étude plus approfondie pourrait montrer s'il existe d'autres solutions, mais cela dépasse le cadre exact sans logarithme ni analyse plus avancée.
\end{correction}

%================================================
\section{Fiche méthodes express}

\begin{rappelbox}
\textbf{1. Simplifier}
\[
\e^a\e^b=\e^{a+b},
\qquad
\frac{\e^a}{\e^b}=\e^{a-b}.
\]

\textbf{2. Résoudre}
\[
\e^A=\e^B \Longleftrightarrow A=B.
\]

\textbf{3. Comparer}
\[
\e^A<\e^B \Longleftrightarrow A<B.
\]

\textbf{4. Dériver}
\[
(\e^x)'=\e^x.
\]

\textbf{5. Étudier un signe}
\[
P(x)\e^x
\]
a le même signe que $P(x)$ car $\e^x>0$.

\textbf{6. Modéliser}
\[
f(t)=C\e^{kt}.
\]
Si $k>0$, croissance. Si $k<0$, décroissance.
\end{rappelbox}

\end{document}

